\documentclass[../../../include/open-logic-section]{subfiles}

\begin{document}

\olfileid{sth}{story}{extensionality}
\olsection{عنصراں نال تعیّن دا اصول}

سبھ توں پہلی گل ایہ اے کہ سیٹ اپنے !!{element}s نال پچھانے جاندے نیں۔
ہور ٹھیک لفظاں وچ:

\begin{axiom}[عنصراں نال تعیّن دا اصول]
جے سیٹ $A$ تے $B$ دے !!{element}s اکّو نیں، تاں $A$ تے $B$ اکّو سیٹ نیں۔
\[
  \lforall[A][\lforall[B][(\lforall[x][(x \in A \liff x \in B)] \lif
  \eq[A][B])]]
\]
\end{axiom}

اسیں ایہ گل \olref[sfr][][]{part} وچ ہر تھاں منّ کے چلے ساں۔ پر ایہ گل
دوبارہ آکھن جوگی اے۔ عنصراں نال تعیّن دا مسلّمہ ایہ بنیادی خیال بیان کردا اے
کہ سیٹ دا تعیّن اوہدے !!{element}s نال ہوندا اے۔ (ایس لحاظ نال سیٹاں نوں
\emph{تصوراں} توں وکھ سمجھ سکدے آں، کیونکہ عین اوہو چیزاں کئی وکھ وکھ
تصوراں تھلّے آ سکدیاں نیں۔)

ایہ اصول کیوں منّیے؟ ایہ کہنا واجب لگدا اے کہ جیہڑا وی عنصراں نال
تعیّن توں انکار کرے، اوہ اوہو جیہی ہر شے چھڈّن دا فیصلہ کردا اے جیہنوں
\emph{سیٹ نظریہ} آکھیا وی جا سکے۔ سیٹ نظریہ عین اوہناں اکٹھاں دا نظریہ اے
جیہناں دا تعیّن اوہناں دے عنصراں نال ہوندا اے۔

پر \olref[sth][][]{part} وچ اصلی اوکڑ ایسے اصول دین دی اے جیہڑے دسّن کہ
\emph{کیہڑے سیٹ موجود نیں}۔ تے پتا لگدا اے کہ ایس سوال دا اکّو سچ مچ
’ظاہر‘ جواب ثابت کیتا جا سکدا اے کہ غلط اے۔

\end{document}
